\chapter{Learning Paradigms: ML, DL and Neurosymbolic AI}
\label{ch:learning-paradigms}
%running head
\chaptermark{Learning Paradigms: ML, DL and NeSy AI}

The previous chapters fixed what a detector sees and what it has to find.
This chapter describes the ways a detector can be built. It starts from
machine learning and deep learning, covers symbolic AI, and then reaches
neurosymbolic AI, where the two meet. It ends with a comparison of the
paradigms and with the question of what an operator needs from an alert.

% ================================================================
% 4.1 Machine learning essentials
% ================================================================
\section{Machine learning essentials}
\label{sec:ml-essentials}

A learning algorithm builds a classifier from examples and not from rules
written by hand. This section covers the part of machine learning that the
benchmark relies on. It describes how the data are labelled, which classical
models exist, what limits their generalization and how class imbalance is
treated.

\subsection{Supervised, unsupervised, semi-supervised learning}
\label{sec:supervised-unsupervised}

The three settings differ in the labels they use. In supervised learning the
model receives pairs of an input and a label and learns to predict the label
of new inputs. Unsupervised learning receives only unlabelled data and looks
for structure in it. Semi-supervised learning uses labelled and unlabelled
data together, so that a large unlabelled set can add to a small labelled
one. The survey of \citet{vanengelen2020semisupervised} places it between
the other two.

Detection shows why the choice matters. An anomaly detector aimed at novel
attacks cannot be trained on the attacks of interest. This is the closed world
problem raised in the introduction. Sommer and Paxson note that one then often
trains the detector with the opposite of what it is supposed to find
\citep{sommer2010outside}. A supervised detector learns the attack classes it has seen. The detectors of
this thesis are supervised. The KDD data are labelled as normal
or as a specific attack \citep{tavallaee2009kddcup99}, and the models of the
benchmark are trained on such labels. One stress scenario reduces the number
of labelled examples available for training. It tests how a supervised
detector behaves under label scarcity.

\subsection{Classical models used in NIDS}
\label{sec:classical-models}

This subsection describes five classical models. How they were used on
\gls{nids} data is the subject of \cref{ch:related-work}.

A decision tree classifies an example by a sequence of tests on its
attributes. \citet{quinlan1986induction} summarizes an approach to building
decision trees from examples and describes one system, ID3, in detail. In our
view a tree is close to a rule set, because each path from the root to a leaf
reads as a condition and a class.

A random forest combines many trees. \citet{breiman2001random} defines it as
a combination of tree predictors in which each tree depends on a random
vector, sampled independently and with the same distribution for all trees.
As trees are added, the generalization error converges to a limit. The limit
depends on the strength of the single trees and on the correlation among
them. Selecting features at random at each split gives error rates that
compare favourably with those of Adaboost, and the method is more robust to
noise.

XGBoost is a tree boosting system. \citet{chen2016xgboost} present it as a
scalable end-to-end system. Their contributions are an algorithm that handles
sparse data, a weighted quantile sketch for approximate tree learning, and
insights on cache access patterns, data compression and sharding. The system
scales beyond billions of examples with far fewer resources than existing
systems.

A support vector machine takes a different route. \citet{cortes1995supportvector}
map the input vectors nonlinearly into a feature space of high dimension. A
linear decision surface is built in that space. The authors extend the method
to training data that cannot be separated without errors.

The nearest neighbor rule assigns to a new point the class of the nearest of
a set of points already classified. \citet{cover1967nearest} show that, in
the large-sample limit, the error of the rule lies between the Bayes error and
twice that value. The bound does not depend on the distribution of the data.

The five models store different things. A tree and a forest keep tests on
features, boosting adds trees one after the other, a support vector machine
keeps a separating surface and the nearest neighbor rule keeps the
classified points themselves. None of them has a place for a domain rule in
its standard form. The conclusion is our own, and it marks the gap that the
neurosymbolic approach addresses.

\subsection{Bias-variance, generalization, regularization}
\label{sec:bias-variance-generalization}

A detector is useful only on traffic it has not seen. \citet{goodfellow2016deep}
call this generalization, the ability to perform well on previously
unobserved inputs. A model underfits when it cannot reach a low error on the
training set. It overfits when the gap between the training error and the
test error is large. The capacity of the model moves it from one case to the
other.

The same authors express the trade-off through the estimator. Bias is the
expected deviation of the estimator from the true value. Variance is how much
the estimator changes from one data sample to another. The mean squared error
is the sum of the squared bias and the variance. Raising the capacity of a
model usually lowers the bias and raises the variance, so the generalization
error follows a U-shaped curve \citep{goodfellow2016deep}.

Regularization covers the changes to a learning algorithm that aim to reduce
the generalization error and not the training error. Weight decay is an
example. It adds to the loss a penalty on the size of the weights
\citep{goodfellow2016deep}.

Two points follow for the benchmark. First, the settings that control
capacity and regularization must be chosen without the test set
(\cref{sec:no-leakage}). Arp et al.
call the opposite practice data snooping, and recommend that test data stay
apart until the final evaluation \citep{arp2022dosdonts}. Second, a fair
comparison needs the same capacity on both sides, which is the
matched-capacity requirement of \cref{sec:matched-capacity}, and the same regularization.
There is also a hypothesis to test. Domain axioms restrict what a model may
learn, and they may therefore act like a regularizer. This is a hypothesis of
the thesis and not a result, and the experiments test it.

\subsection{Handling class imbalance}
\label{sec:imbalance-handling}

The effect of imbalance on evaluation was covered in \cref{sec:imbalance-leakage-bias}.
Here the question is how a model can be trained on classes of very different
size. Three families of methods exist.

The first changes the data. \citet{chawla2002smote} propose SMOTE, which
builds artificial examples of the minority class and does not only remove
majority examples. For a minority example, SMOTE selects some of its nearest
neighbors, takes the difference vector to each of them, multiplies it by a
random number between 0 and 1 and adds it to the example. The new point lies
on the segment between the two. The authors report that combining it with
undersampling of the majority class gives better results than undersampling
alone. The same combination also gives better results than a change of the
loss ratio in Ripper or of the class priors in Naive Bayes. The classifiers of
the evaluation are C4.5, Ripper and Naive Bayes, and performance is measured
with ROC curves and the area under the curve.

SMOTE-NC extends the method to data with nominal and continuous features. It
adds a penalty to the Euclidean distance when nominal values differ. The
penalty is the median of the standard deviations of the continuous features
of the minority class. A new example takes the nominal value held by the
majority of its $k$ nearest neighbors \citep{chawla2002smote}. Which of the benchmark datasets
mix the two kinds of feature is described with the data preparation (\cref{sec:impl-data}).

Two cautions are ours. A synthetic point is an interpolation and not an
observed flow, so nothing guarantees that it respects the domain rules of
\cref{ch:attacks-ddos}. And resampling must be fitted on the training split
only. A synthetic point built from a test example would let test information
reach training, which is the data snooping pitfall again
\citep{arp2022dosdonts}.

The second family changes the loss and leaves the data alone, so that errors
on the minority class cost more. \citet{chawla2002smote} compare SMOTE with
such an adjustment, the loss ratio in Ripper. For deep networks,
\citet{johnson2019survey} group the methods for class imbalance into
data-level, algorithm-level and hybrid ones. Algorithm-level methods are
commonly implemented with a weight or cost schema, and the authors find that
cost-sensitive learning remains applicable to deep learning. How the weights
are set for the neural models is fixed in the benchmark design.

The third family is the focal loss. \citet{lin2017focal} propose it for
one-stage object detectors, where the many easy background examples overwhelm
training. It changes the cross-entropy loss so that well-classified examples
count less and training concentrates on the hard ones. The setting of the
paper is object detection and not traffic analysis. Using the focal loss on
flows is a choice that the benchmark has to justify.

Resampling acts on the data, while class weighting and the focal loss act on
the training objective. The \gls{ltn} approach of this thesis also acts
through the training objective, as \cref{ch:ltn} shows.

% ================================================================
% 4.2 Deep learning
% ================================================================
\section{Deep learning}
\label{sec:deep-learning}

Deep learning builds a model by composing many layers. A feedforward network
approximates a function as a chain of simpler functions, and its depth is the
number of layers in the chain \citep{goodfellow2016deep}. This section covers
the networks of the deep baselines, how they are trained, which architectures
exist for sequential and unlabelled data, and where the paradigm is strong and
weak.

\subsection{Multilayer perceptrons and their training}
\label{sec:mlp-training}

A multilayer perceptron is a feedforward network. Each layer applies an affine
transformation followed by a nonlinear activation function. The rectified
linear unit, which returns the larger of zero and its input, is a common
choice. The training data do not specify the output of the hidden layers, and
the learning algorithm decides how to use them \citep{goodfellow2016deep}.

The network is trained with an iterative, gradient-based optimizer that lowers
a cost function. Unlike convex methods, this gives no guarantee of reaching a
global minimum. In most modern networks the cost is the cross-entropy between
the training data and the model distribution \citep{goodfellow2016deep}.
Backpropagation computes the gradient that the optimizer needs. It is a
method for computing the gradient and not a learning algorithm. It applies the
chain rule through the computation graph and stores intermediate values, so
its cost grows linearly with the number of edges of the graph
\citep{goodfellow2016deep}.

Computing the gradient on the whole training set is expensive. Stochastic
gradient descent estimates it on random minibatches. Because the noise of the
estimate never disappears, the learning rate has to decrease during training
\citep{goodfellow2016deep}. Adam adapts the step of each parameter. Kingma
and Ba base it on adaptive estimates of lower-order moments and describe it
as simple, efficient and light on memory, suited to noisy or sparse gradients
\citep{kingma2015adam}. The textbook treatment adds that it is fairly robust
to the choice of its hyperparameters \citep{goodfellow2016deep}.

Two techniques limit overfitting. Dropout removes units and their connections
at random during training. This keeps units from co-adapting too much, and
training samples from an exponentially large set of thinned networks. At test
time a single unthinned network with smaller weights approximates the average
of their predictions \citep{srivastava2014dropout}. Early stopping watches the error on a
validation set during training and keeps the best parameters seen, so the
number of training steps works as a control on capacity
\citep{goodfellow2016deep}. It needs a validation set separate from the test
set, which is one reason why the benchmark keeps three splits.

\subsection{Deep architectures for traffic}
\label{sec:dl-architectures}

This subsection describes four families of architecture. \Cref{ch:related-work}
reviews how they have been applied to network traffic.

A convolutional layer scans the input with a small kernel and applies the same
linear transformation at every position. It relies on three ideas. Interactions
are sparse, the weights of the kernel are shared across positions, and the
representation is equivariant to shifts of the input. Convolution also applies
to one-dimensional data such as a time series \citep{goodfellow2016deep}. The
features of a flow record are not a time series. They have no natural
neighbourhood. A one-dimensional convolution over a feature vector therefore
assumes an order that comes from the layout of the features, a limit we note
for the 1D-CNN baseline.

Recurrent networks read a sequence step by step. The long short-term memory
unit is the standard gated variant. Greff et al. evaluate eight variants on
speech recognition, handwriting recognition and music modelling. They find that
none of the variants improves substantially on the standard unit and that the
forget gate and the output activation function are its most critical
components \citep{greff2015lstm}. The gated recurrent unit is a second gated
design. Chung et al. compare gated units with plain tanh units on music and
speech modelling. Both kinds of gated unit do better, and the gated recurrent
unit performs comparably to the long short-term memory unit
\citep{chung2014empirical}. The benchmark plan lists an MLP and a 1D-CNN as
deep baselines and an LSTM as an optional extension.

An autoencoder has an encoder that maps the input to a code and a decoder that
rebuilds the input from the code. It is kept from copying its input exactly.
When the code is smaller than the input the autoencoder is undercomplete, and
with a linear decoder and a squared error loss it learns the same subspace as
principal component analysis. A denoising autoencoder receives a corrupted
input and rebuilds the clean one \citep{goodfellow2016deep}. The linear case
connects to the PCA detectors of \cref{sec:statistical-anomaly-detection}. For
detection, an autoencoder can be trained on normal traffic alone, and a large
reconstruction error is then read as a deviation \citep{mirsky2018kitsune}.
We state this as an option,
and \cref{sec:deep-nids} reviews a detector built this way.

The transformer is the newest of the four. Vaswani et al. propose an
architecture that relies on attention alone and drops recurrence and
convolution. They report that it parallelizes better and trains faster than
the recurrent and convolutional models of the time on machine translation
\citep{vaswani2017attention}.

\subsection{Strengths and weaknesses of deep learning}
\label{sec:dl-strengths-weaknesses}

Two strengths stand out. The first is representation learning. The hidden
layers of a network are not specified by the data, so the model builds its own
features \citep{goodfellow2016deep}. The lower-dimensional codes of an
autoencoder can improve classification and reduce cost
\citep{goodfellow2016deep}. The second is scale. Stochastic gradient descent
makes training on large datasets practical, because the computation time of an
update does not grow with the number of training examples
\citep{goodfellow2016deep}. Attention-based models
add parallel training \citep{vaswani2017attention}.

The weaknesses matter more for this thesis, since each one points at a
condition of the benchmark.

\emph{Data hunger.} Srivastava et al. state that overfitting is a serious
problem in networks with a large number of parameters
\citep{srivastava2014dropout}. A model with many parameters needs many
labelled examples to avoid it. We expect the effect to show under label
scarcity, and the corresponding stress scenario tests this expectation.

\emph{Opacity.} A deep network returns a score and no reason. Sommer and Paxson
name the semantic gap between a detection and a report that an operator can
act on \citep{sommer2010outside}. We judge that a network with many layers
widens that gap.

\emph{Brittleness under shift.} Arp et al. describe spurious correlations,
where a model learns artifacts that have nothing to do with the security
problem, and lab-only evaluation, which ignores that data change over time
\citep{arp2022dosdonts}. A network fitted on one collection can carry the
artifacts of that collection to the next. Cross-dataset transfer tests this in
the benchmark.

\emph{No constraint guarantees.} The loss measures how well the predictions
fit the labels. Nothing in it forbids a prediction that breaks a domain rule,
such as one of the candidate axioms of \cref{ch:attacks-ddos}. The domain-rule
violation rate of \cref{sec:contributions} measures how often this happens.

None of the four weaknesses is specific to network traffic, and none is
proved here. Apart from the cited findings, the four readings are ours. They
state what the benchmark has to check, and they are the reason why a model
trained with explicit domain knowledge is worth comparing with a purely
neural one.


% ================================================================
% 4.3 Symbolic AI
% ================================================================
\section{Symbolic AI}
\label{sec:symbolic-ai}

Symbolic AI represents knowledge in a formal and declarative form, based on
rules and logic \citep{colelough2025neurosymbolic2024}. Colelough and Regli
name DENDRAL and MYCIN among the early systems of this kind. In this thesis a
symbolic system is one whose knowledge is stated as formulas that a program
can reason over.

\subsection{Propositional and first-order logic}
\label{sec:logic-inference}

Propositional logic builds formulas from atoms, each true or false, and from
connectives such as and, or, not and implication. First-order logic adds
objects, predicates over objects and quantifiers. Russell and Norvig give
logical agents, first-order logic and inference in first-order logic a
chapter each \citep{russell2020aima}.

A Horn clause has at most one positive literal. Read as a rule, it says that
a conjunction of atoms, the body, implies a single atom, the head. The SYN
flood axiom of \cref{tab:candidate-axioms} has this form: a window with a SYN
excess and incomplete connections implies DoS, the KDD family that includes
DDoS. Two procedures apply such
rules. Forward chaining starts from known facts and derives every head whose
body holds. Backward chaining starts from a query and looks for rules whose
head matches it.

Inference is exact. If the facts and the rules are true, the derived
conclusion is true. The rules of \cref{ch:attacks-ddos} do not meet this
condition yet. Their predicates, such as a high rate, need a threshold before
they are true or false, and the chapter left the thresholds open. This gap,
as we see it, is the reason why \cref{ch:ltn} replaces truth values by degrees
of truth.

\subsection{Expert systems, ontologies and knowledge graphs}
\label{sec:rules-ontologies-kg}

A rule-based system keeps its rules in a knowledge base and applies them with
a separate inference engine. We regard the signature rules of
\cref{ch:foundations} as a narrow case of such a system. A Snort rule combines
options by a logical AND \citep{roesch1999snort}.

An ontology is a formal representation of what the terms of a domain mean
\citep{hogan2021knowledge}. The Web Ontology Language is among the most
popular languages for writing one, and it is a recommendation of the W3C
\citep{hogan2021knowledge}. Description logics are a family of knowledge
representation languages that are widely used to model ontologies, and they
are one of the main foundations of the Web Ontology Language
\citep{krotzsch2012primer}.

Hogan et al. define a knowledge graph as a graph of data meant to accumulate
and convey knowledge of the real world, where nodes stand for entities and
edges for relations between them \citep{hogan2021knowledge}. In network
security the nodes could be hosts, services or ATT\&CK techniques. The
thesis uses ATT\&CK only as a label space. How a graph can enter a model is
discussed later in this chapter.

\subsection{Answer set and probabilistic logic programming}
\label{sec:asp-plp}

Answer Set Programming describes a problem by rules and asks for the sets of
atoms that satisfy them. Brewka et al. present rules with a head and a body,
default negation for incomplete information, and answer sets as the
solutions \citep{brewka2011answer}. They name search and optimization
problems as the target and stress the declarative style of the approach.

Probabilistic logic programming extends languages such as Prolog with
probabilities. De Raedt and Kimmig survey the concepts behind these languages
and their execution mechanisms \citep{deraedt2015probabilistic}. In our own
summary, a query to such a program returns a probability and not only true
or false. Both families return later in this chapter as the symbolic half of
neurosymbolic systems.

\subsection{Strengths and weaknesses of symbolic AI}
\label{sec:symbolic-strengths-weaknesses}

Garcez and Lamb list the strengths of symbolic AI. It reasons in a
sophisticated way, and it can give a trace of the steps that led to an
outcome \citep{garcez2020thirdwave}. They also note that reasoning is often
meant to be precise and not approximate. A rule base needs no training
examples, because the knowledge comes from the expert. This remark on data
efficiency is ours, and so is the observation that a rule can be checked
against a case, since every conclusion has a derivation.

The same authors name brittleness and computational complexity among the
limits of symbolic systems \citep{garcez2020thirdwave}. The thesis uses the
term knowledge acquisition bottleneck for the cost of obtaining the rules, a
term that Kautz uses in his account of expert systems \citep{kautz2022third}.

A third limit is the lack of perception. A symbolic system reasons over
symbols and cannot read raw input. A flow record has to become predicates
first, and the thresholds that do this are themselves a choice, as
\cref{sec:logic-inference} showed. This limit is our reading and not a claim
of the sources.

Set against \cref{sec:dl-strengths-weaknesses}, the two paradigms fail in
opposite places. A neural model needs data, hides its reasoning and ignores
rules. A symbolic one needs no data and shows its reasoning, but it cannot
read raw input. This contrast is ours, and the next section builds on it.

% ================================================================
% 4.4 Neurosymbolic AI
% ================================================================
\section{Neurosymbolic AI}
\label{sec:neurosymbolic-ai}

\subsection{Why combine the two}
\label{sec:complementarity}

The usual argument for combining the two paradigms draws on Kahneman. System 1
is fast and intuitive and System 2 is slow and deliberate. Garcez and Lamb
say that neural networks and symbolic reasoning can in principle be matched
with the two systems, respectively, and they report that Kahneman said a
symbolic System 2 layer is needed on top of deep learning
\citep{garcez2020thirdwave}. Colelough and Regli use the same framework to
justify the combination \citep{colelough2025neurosymbolic2024}.

For traffic, we map the two systems as follows. The neural part reads the
numbers of a flow and the symbolic part holds the rules of
\cref{ch:attacks-ddos}.

The analogy has a limit. It says what each part does well and not how the two
are joined. It is a motivation and not evidence that a combined system
detects better. The benchmark tests that question.

\subsection{Kautz's taxonomy of architectures}
\label{sec:kautz-taxonomy}

Kautz describes six designs for joining neural and symbolic parts, and he
gives each a short notation \citep{kautz2022third}. \Cref{tab:kautz-types}
lists them with that notation and a paraphrase of his descriptions. Garcez and
Lamb also report six types, one of which is capable of true symbolic reasoning
inside a neural engine \citep{garcez2020thirdwave}.

\begin{table}[htbp]
\centering
\footnotesize
\setlength{\tabcolsep}{3pt}
\caption{Kautz's six types, with the notation of the original and a
paraphrase of its descriptions \citep{kautz2022third}.}
\label{tab:kautz-types}
\begin{tabular}{@{}p{0.8cm}p{3.4cm}p{8.3cm}@{}}
\toprule
Type & Notation & Description \\
\midrule
1 & Symbolic Neuro symbolic
  & Symbolic input goes through a neural network and the output is symbolic. \\
2 & $\mathrm{Symbolic}[\mathrm{Neuro}]$
  & A symbolic problem solver uses a neural network as a pattern-recognition subroutine. \\
3 & $\mathrm{Neuro}\,|\,\mathrm{Symbolic}$
  & A neural network turns non-symbolic input into a symbolic form on which a symbolic part reasons. \\
4 & $\mathrm{Neuro}:\mathrm{Symbolic}\rightarrow\mathrm{Neuro}$
  & A training regime based on symbolic rules shapes the network. \\
5 & $\mathrm{Neuro}_{\mathrm{Symbolic}}$
  & Symbolic rules become templates for structures inside the network. \\
6 & $\mathrm{Neuro}[\mathrm{Symbolic}]$
  & A symbolic reasoning engine is embedded in a neural system, which calls it when needed. \\
\bottomrule
\end{tabular}
\end{table}

The accounts do not agree on every type. \citet{wan2024cognitive} list five
paradigms and leave out the first. They describe the type with a symbolic
training regime as rules that guide learning, with the knowledge compiled into
the structure of the network. They describe the type with templates as rules
mapped onto embeddings that act as soft constraints or regularizers. For the
type with an embedded symbolic part they speak of explainability and
robustness taken from symbolic reasoning, with attention to symbolic
information. Their wording differs from that of \cref{tab:kautz-types}.

Kautz names logic tensor networks as an example of the fifth type
\citep{kautz2022third}, and \citet{wan2024cognitive} describe that type as
soft constraints or regularizers. The models of this thesis are logic tensor
networks, and a loss that contains the formulas of a knowledge base is such a
constraint. We therefore read them as the fifth type. This is our judgment,
and \cref{ch:ltn} keeps it with a reservation. The families of the next
subsection differ in how the neural and the symbolic components are joined,
which is what a type records.



\subsection{Integration families}
\label{sec:integration-families}

The systems below differ in where the logic acts. It can act on the loss, inside
a logic program, on a graph, on the structure of the network or on the
output. The order of the subsections is not a ranking.

\subsubsection{Logic as a loss}
\label{sec:logic-as-loss}

Logic Tensor Networks use Real Logic, a many-valued and end-to-end
differentiable first-order logic. Badreddine et al. apply it to tasks such as
multi-label classification, relational learning, query answering and
semi-supervised learning, with an implementation in TensorFlow 2
\citep{badreddine2022ltn}. The semantic loss of Xu et al. is a loss that
connects the output vector of a network with logical constraints, and the
authors show gains in semi-supervised classification and in the prediction of
structured objects \citep{xu2018semantic}. DL2 lets the user state domain
knowledge as declarative constraints that are enforced during training and
translates them into loss terms for gradient-based optimization. The same
system can query a trained model, and the authors claim a higher
expressiveness than earlier approaches \citep{fischer2019dl2}.

What the three share, in our summary, is that a formula becomes a
differentiable term of the training objective. This places the family on the
same side as the loss-level treatments of imbalance in
\cref{sec:imbalance-handling}. We also note that a penalty makes a violation
costly but does not forbid it, so the fourth weakness of
\cref{sec:dl-strengths-weaknesses} is not removed. \Cref{ch:ltn}
states why the thesis takes this family.

\subsubsection{Neural predicates in a logic program}
\label{sec:neural-predicates}

In DeepProbLog a neural predicate takes its probabilities from the softmax
output of a network, so that the logic program treats the network as a source
of probabilistic facts. Manhaeve et al. claim to be the first to integrate
general-purpose neural networks with expressive probabilistic-logical
modelling and reasoning \citep{manhaeve2018deepproblog}. DeepStochLog works
with stochastic definite clause grammars, a kind of stochastic logic program,
and adds neural grammar rules. Its authors report that inference and learning
scale much better than in neural probabilistic logic programs
\citep{winters2022deepstochlog}. NeurASP reads the outputs of a network as
probability distributions over facts of an answer set program. It can use
pre-trained networks for symbolic reasoning, and it can use the rules to
improve the training of the network \citep{yang2020neurasp}.

For traffic, we picture a network that scores predicates such as a high
rate, and a program that derives the attack class from them. The three systems were presented on other tasks, and the
sources read here do not report network traffic.

\subsubsection{Knowledge-graph and embedding-based integration}
\label{sec:kg-integration}

DeLong et al. survey neurosymbolic reasoning over knowledge graphs and group
the methods in three categories. Logically informed embeddings use symbolic
inference to augment the graph before a neural embedding. Embeddings with
logical constraints impose rules on the embedding space, on the loss or
through separate logic embeddings. Rule learning approaches learn the rules
themselves \citep{delong2024kgsurvey}. In security a graph could link hosts,
services and ATT\&CK techniques. The thesis does not build such a graph.

\subsubsection{Neural module networks and program synthesis}
\label{sec:neural-modules}

Andreas et al. assemble jointly trained neural modules into a network for
each question about an image. They parse the question into a linguistic
structure and use it to build the modular network from reusable components
\citep{andreas2015neuralmodule}. The structure of the computation therefore
comes from the input. A flow record carries no such structure, so we see no
natural place for this family in the benchmark.

Program synthesis is a neighbouring topic. Chaudhuri et al. survey
neurosymbolic programming. There the learned functions are programs that can
combine neural modules with symbolic primitives, and they are found by symbolic
search together with gradient-based optimization
\citep{chaudhuri2021neurosymbolic}. The benchmark starts from rules written in
advance and does not search for a program, so we leave this family out.

\subsubsection{Post-hoc symbolic verification of neural outputs}
\label{sec:post-hoc-verification}

Here the model is left as it is and checked afterwards. Katz et al. verify
properties of trained networks with an extension of the simplex method to
ReLU activations, and they report that it proves properties of networks an
order of magnitude larger than the largest verified with existing methods
\citep{katz2017reluplex}. Renkhoff et al.
survey verification, validation and testing of neurosymbolic systems
\citep{renkhoff2024survey}. The simplest form of the idea is to test every
prediction against a rule and to flag the ones that break it. The domain-rule
violation rate of \cref{sec:contributions} is such a check, applied to the
outputs of every model, in our interpretation. The check measures a model and
does not change how it is trained.

\subsection{Comparison of the families}
\label{sec:integration-comparison}

\Cref{tab:integration-systems} puts the systems side by side. It lists only
what the sources read here state.

\begin{table}[htbp]
\centering
\footnotesize
\setlength{\tabcolsep}{3pt}
\caption{Integration systems, as described in the sources read for this section.}
\label{tab:integration-systems}
\begin{tabular}{@{}p{2.3cm}p{5.2cm}p{5.0cm}@{}}
\toprule
System & Symbolic part & Role of the network \\
\midrule
LTN & Real Logic, many-valued differentiable first-order logic
  & Trained end to end together with the logic \\
Semantic loss & Logical constraints
  & Its output vector is tied to the constraints by the loss \\
DL2 & Declarative constraints, also usable as queries
  & Trained under a loss derived from the constraints \\
DeepProbLog & Probabilistic logic program
  & Neural predicates give the probabilities \\
DeepStochLog & Stochastic definite clause grammars
  & Neural grammar rules \\
NeurASP & Answer set program
  & Outputs act as distributions over facts \\
\bottomrule
\end{tabular}
\end{table}

The sources do not use one scale, so the dimensions below are compared
only where a source speaks.

\emph{Expressiveness.} The authors of DL2 claim a higher expressiveness than
earlier approaches, and the authors of DeepProbLog claim an expressive
probabilistic-logical component. Each claim is the claim of its authors and
rests on their own benchmarks. We do not rank the systems.

\emph{Scalability.} The authors of DeepStochLog report that stochastic logic
programs scale much better than neural probabilistic logic programs. No source
read here compares the loss-based systems with the program-based ones.

\emph{Training cost.} No source read here reports it for network traffic. We
expect a logic loss to add the cost of evaluating the formulas at every
step, and a logic program to add the cost of inference in that program. The
benchmark measures training time directly.

\emph{Explainability.} A logic program can show the derivation of a
conclusion, in line with the trace that \citet{garcez2020thirdwave} count
among the strengths of symbolic AI. A model trained with a logic loss is an
ordinary network once training ends, unless the logic is kept at inference.
This is our reading, and \cref{ch:ltn} checks it for LTN.

\emph{Tooling maturity.} The only statement read here is that Logic Tensor
Networks have an implementation in TensorFlow 2 \citep{badreddine2022ltn}.
The maturity of the other tools has not been compared.

The choice of the thesis rests mainly on the data. The inputs are tabular flow
features and the knowledge is a set of rules over predicates of flows
(\cref{ch:attacks-ddos}). A logic loss fits a training pipeline like the one
of the deep baselines, and we expect the cost at inference to stay close to
that of the baseline. This is a hypothesis of the thesis, and the benchmark
tests it.


% ================================================================
% 4.5 Comparison of the paradigms
% ================================================================
\section{Comparison of the paradigms}
\label{sec:paradigm-comparison}

\Cref{tab:paradigm-comparison} compares the four paradigms on seven
dimensions. Every cell is an expectation of ours, drawn from
the sections above, and none is a result. The last column describes a logic
loss, the family of the thesis. Other families differ.

\begin{table}[htbp]
\centering
\footnotesize
\setlength{\tabcolsep}{3pt}
\caption{Our expectations for the four paradigms, not results. The benchmark
tests the rows on data, cost and maintenance.}
\label{tab:paradigm-comparison}
\begin{tabular}{@{}p{2.4cm}p{2.5cm}p{2.5cm}p{2.5cm}p{2.5cm}@{}}
\toprule
 & Classical ML & Deep learning & Symbolic & NeSy, logic loss \\
\midrule
Data efficiency & Needs labels & Needs many labels & Needs rules, no labels & Labels and rules \\
Interpretability & A tree is close to rules, other models less & Low
  & Derivation of each conclusion & An ordinary network unless the logic is kept \\
Constraint compliance & None & None & Exact & Soft, not guaranteed \\
Training cost & Expected lower & Expected higher & Writing the rules & Deep cost plus the formulas \\
Inference cost & Expected low & A forward pass & Inference in the program & Expected as deep, a hypothesis \\
Maintenance cost & Retrain under shift & Retrain under shift & Edit the rules by hand & Retrain and edit the rules \\
Maturity & Models cited date from 1967 to 2016 & Works cited date from 2014 to 2017
  & Early systems such as DENDRAL and MYCIN & Logic-loss systems cited date from 2018 to 2022 \\
\bottomrule
\end{tabular}
\end{table}

Three points stand out to us. No paradigm wins on every row. The cost of a
symbolic system is the work of writing the rules and not the machine time.
The two learning paradigms need retraining when the data change, and
\citet{arp2022dosdonts} note that data change over time.

% ================================================================
% 4.6 Explainable AI in this context
% ================================================================
\section{Explainable AI in this context}
\label{sec:xai}

\subsection{Post-hoc and intrinsic explanations}
\label{sec:posthoc-vs-intrinsic}

Post-hoc methods explain a model that is already trained. LIME explains the
prediction of any classifier by learning an interpretable model locally
around the prediction \citep{ribeiro2016lime}. SHAP introduces a class of
additive feature importance measures that unifies six earlier methods
\citep{lundberg2017shap}. Both attribute a prediction to input features. For
a flow record the answer is a list of features, such as a packet rate or a
count of flags, with a weight each.

Rudin argues that in high-stakes decisions one should not explain black box
models after the fact and should use inherently interpretable models instead
\citep{rudin2019stop}. A structured explanation is of the second kind. In a
rule-based system it is the derivation that led to the conclusion, which
Garcez and Lamb count among the strengths of symbolic AI
\citep{garcez2020thirdwave}. Kalutharage et al. pair neurosymbolic learning
with explanation driven by domain knowledge for the detection of attacks on
IoT networks \citep{kalutharage2025neurosymbolic}.

A model trained with a logic loss does not give a derivation by itself, as
\cref{sec:integration-comparison} noted. It may need a post-hoc method as
well, unless the rules are evaluated again at inference. We add this point
ourselves.

\subsection{What an operator actually needs in an alert}
\label{sec:operator-needs}

The semantic gap of Sommer and Paxson returns here from the side of the
operator. The same authors stress the cost of errors, since a false positive
takes the time of an analyst \citep{sommer2010outside}. We read four
needs from this and from the defences of \cref{ch:attacks-ddos}. The alert
should say what was seen, which flow, window or source. It should say why, in
terms of features or rules. It should say how sure the detector is. It should
name the traffic in fields that a filter can use, since filtering needs a
description of the attack traffic. This list is ours.

Besides the engineering effort, RQ4 asks what the approach returns in alerts
that a human analyst can interpret and audit. A feature weight meets the second need only in part,
because it does not say which condition of an attack held. A rule that fired
names one, although beside a learned model it agrees with the decision and does
not cause it. Whether this difference matters to an analyst is an empirical
question, and the usability measurements of \cref{sec:usability-dimension}
address it.

% ================================================================
% 4.7 Summary
% ================================================================
\section{Summary}
\label{sec:ch4-summary}

The machine learning and deep learning models of the benchmark learn from
labelled data, and they offer no place for domain rules. Symbolic AI holds rules and gives derivations, but
it cannot read raw input and its rules are costly to obtain. Neurosymbolic AI
joins the two, and the families differ in where the logic acts: on the loss,
inside a logic program, on a graph, on the structure of the network or on the
output. Kautz's taxonomy records how the parts are joined, and the thesis
reads its models as the fifth type, which is its own judgment. The thesis
uses a logic loss, which fits tabular flow features and rules over predicates
of flows. The comparison
of the paradigms is a set of expectations, and the benchmark tests those on
data, cost and maintenance. Explanations are a further point on which the
paradigms differ for an operator, and RQ4 asks about it.